\documentclass[10pt,a4paper]{amsart}
\usepackage[margin=19mm]{geometry}
\usepackage{amsmath,amssymb,mathtools}
\usepackage[scr=rsfs,cal=euler]{mathalfa}
\usepackage{xfrac}
\usepackage[unicode,hidelinks,bookmarks=false]{hyperref}
\newcommand{\ie}{i.e.\ }
\newcommand{\cf}{cf.\ }
\newcommand{\resp}{respectively\ }
\newcommand{\Subsection}{\subsection}
\providecommand{\nobreakdash}{\nobreak\hbox{-}}
\setlength{\emergencystretch}{2em}
\multlinegap0pt
\usepackage[shortlabels]{enumitem}
\usepackage{amssymb}
\usepackage{xcolor}
\usepackage[normalem]{ulem}
\newtheorem{proposition}{Proposition}
\newtheorem{lemma}{Lemma}
\newtheorem{remark}{Remark}
\newcommand{\N}{\mathbb{N}}
\newcommand{\bes}{\begin{equation*}}
\newcommand{\ees}{\end{equation*}}
\newcommand{\ep}{\varepsilon}
\title{On the complete separation of unique $\ell_{1}$ spreading models and the Lebesgue property of Banach spaces}
\author{Harrison Gaebler, Pavlos Motakis, Bunyamin Sari}
\date{Source: arXiv:2402.14687v2 (CC BY 4.0)}
\begin{document}
\maketitle

\noindent\emph{Source and licence.} On the complete separation of unique $\ell_{1}$ spreading models and the Lebesgue property of Banach spaces. Version arXiv:2402.14687v2 (CC BY 4.0), distributed by arXiv under the Creative Commons Attribution 4.0 International licence, http://creativecommons.org/licenses/by/4.0/, as recorded in the arXiv licence metadata for this exact version and shown on the abstract page https://arxiv.org/abs/2402.14687v2. Full licence notice: \texttt{licenses/arxiv-2402.14687-CC-BY-4.0.txt}.\par\smallskip

\noindent\textit{Editorial scope.} Counted unit: the article's proposition c0\_type\_estimate (source0.tex lines 362--366). The article's complete original proof of it is delivered with it (source0.tex lines 367--403, one \texttt{proof} environment of 37 lines). No mathematics is written, repaired or re-derived: the statement and the proof are the article's own lines, in the article's own notation, and no retained line is substituted or re-flowed.  Retained uncounted as the source-local context the proof needs: lines 133--135; lines 186--190; lines 339--341; lines 354--358.  Nothing was omitted from any retained span, so the package carries no omission record.  No retained line contains a citation, so the wrapper carries no bibliography and no bibliography entry.  No retained line contains a figure environment, an image-inclusion command or drawing code, so no image is delivered and the package declares no asset dependency.  Editorial formatting only, in addition: an AMS \texttt{amsart} preamble that restates the article's own declarations and macros used by the retained lines, namely the environments \texttt{proposition}, \texttt{lemma}, \texttt{remark} on separate counters, together with the standard packages those lines need.  The retained environments print in this document as Proposition 1, Lemma 1, Remark 1 in order of appearance, whereas the article prints the same items with the numbering of its own full document.  The complete native source of the article as distributed in the arXiv e-print for this exact version is bundled unchanged as \texttt{sources/154976/source0.tex}.
\begin{proposition}{\label{basic-scc-maximal}}
Let $n\in\N$. Then, for every infinite subset $M\subset\N$ and for every $F\subset M$ such that $F\in S_{n}$ is maximal with respect to inclusion, there exist real numbers $(c_{i})_{i\in F}$ so that the vector $x=\sum_{i\in F}c_{i}e_{i}\in c_{00}$ is an $(n,3/\min(F))$-basic scc.
\end{proposition}

Fix a pair of strictly increasing sequences $(m_{j})$ and $(n_{j})$ of positive integers such that $m_{1}=100$, $n_{1}=1$,
\begin{enumerate}[label=(\arabic*)]
\item\label{condition m_j} $m_{j+1}\geq j^2m_j^2$, for every $j\in\N$,
\item\label{condition n_j} $n_{j+1}>\frac{4}{\log(100)}n_{j}\log(m_{j+1})$.
\end{enumerate}

\begin{lemma}{\label{aux_lemma}}
Let $f=\frac{1}{m_{j_{1}}\cdots m_{j_{l}}}\sum_{i=1}^{d}f_{i}\in W^{\text{aux},\delta}$. Then, $\{ i \mid f_{i}\notin W^{\text{aux},\delta}_{0}\}=E_{0}\cup E_{1}\cup\ldots\cup E_{l}$ where $\mathrm{card}(E_0)\leq 1$ and, for each $1\leq z\leq l$, $(f_{i})_{i\in E_{z}}$ is $S_{n_{j_{1}}+\cdots+n_{j_{z}}}\ast\mathcal{A}_{3}$-admissible and $\max_{i\in E_{z}}d_{T}(m(f_{i}),\lambda_{z})\leq\delta$ for some $\lambda_{z}\in\overline{\mathcal{D}}$.
\end{lemma}

\begin{remark}
\label{weight-to-admissibility condition}
Let $j\in\N$ and $f\in W$ such  $m(f)<m_j$ and $w(f) < m_j^4$, then, if $\vec w(f)= (m_{j_1},\ldots,m_{j_\ell})$, $n_{j_1}+\cdots +n_{j_\ell} < n_j$. Indeed, $m_{1}^{l}\leq w(f) <m_{j}^{4}$, and thus, $l<4\log(m_{j})/\log(100)$. By condition \ref{condition n_j} on page \pageref{condition n_j},
\[n_{j_{1}}+\cdots+n_{j_{l}}\leq ln_{j-1}< n_{j-1}4\frac{\log(m_{j})}{\log(100)}\leq n_{j}.\]
\end{remark}

\begin{proposition}{\label{c0_type_estimate}}
Let $N\in\N\cup\{0\}$ and choose $(t_{i})_{i=1}^{2^{N}}$ such that $d(m_{t_{i}},m_{t_{i'}})\geq 2^{-N}$ for each $i\neq i'$. Then, choose $\ep_{i}\in\left(0,\frac{1}{3m_{t_{i}}2^{N}}\right)$ and define a $(n_{t_{i}},\ep_{i})$-basic scc by
\bes x_{i}=m_{t_{i}}\sum_{r\in F_{i}}c_{r}^{i}e_{i_{r}} \ees
for each $1\leq i\leq 2^{N}$. It follows that $\left\vert f\left(\sum_{i\in\Omega}a_{i}x_{i}\right)\right\vert\leq \left(1+\frac{4}{\sqrt{w(f)}}\right)\max_{i\in\Omega}|a_{i}|$ for every subset $\Omega\subset\{1,\ldots,2^{N}\}$, for every choice $(a_{i})_{i\in\Omega}$ of scalars, and for every auxiliary weighted functional $f=\frac{1}{m_{j_{1}}\cdots m_{j_{l}}}\sum_{q=1}^{d}f_{q}\in W^{\text{aux},\delta}$ where $\delta<2^{-(N+1)}$.
\end{proposition}

\begin{proof}
Let $\Omega\subset\{1,\ldots,2^{N}\}$ be given. We will proceed by induction on the level $k$ norming sets $W^{\text{aux},\delta}_{k}$ and the inductive hypothesis will include that $\left\vert f\left(\sum_{i\in A_{f}(\Omega)}a_{i}x_{i}\right)\right\vert<\frac{4\gamma}{\sqrt{w(f)}}$ if $f\notin W_{0}^{\text{aux},\delta}$ where $\gamma=\max_{i\in\Omega}|a_{i}|$ and where $A_{f}(\Omega) :=\{i\in\Omega \mid m_{t_{i}}\neq m(f)\}$. The base case is the level $k=1$ norming set and we first let $f\in W^{\text{aux},\delta}_{0}$ so that   
\bes \left\vert f\left(\sum_{i\in\Omega}a_{i}x_{i}\right)\right\vert = \left\vert e^{*}_{i_{0}}\left(\sum_{i\in\Omega}a_{i}m_{t_{i}}\sum_{r\in F_{i}}c^{i}_{r}e_{r} \right)\right\vert \leq \max_{i\in\Omega}\ep_{i}m_{t_{i}}|a_{i}| < \max_{i\in\Omega}|a_{i}| \ees
for some $i_{0}\in\N$. It therefore remains to check functionals of the form
\bes f=\frac{1}{m_{j_{1}}\cdots m_{j_{l}}}\sum_{q=1}^{d}e^{*}_{i_{q}}\in W^{\text{aux},\delta}_{1}\setminus W^{\text{aux},\delta}_{0} \ees
and we write $A_{f}(\Omega)=\{ i\in\Omega \mid m_{t_{i}}\neq m(f)\}=A_{1}\cup A_{2}\cup A_{3}$ where
\begin{align*} A_{1}&=\{ i\in\Omega \mid m_{t_{i}}<m(f)\} \\ A_{2}&=\{ i\in\Omega \mid m(f)<m_{t_{i}}\text{ and } n_{j_{1}}+\cdots+n_{j_{l}}<n_{t_{i}}\} \\ A_{3}&=\{ i\in\Omega \mid m(f)<m_{t_{i}} \text{ and } n_{j_{1}}+\cdots+n_{j_{l}}\geq n_{t_{i}}\} \end{align*}
so that $\Omega\setminus A_{f}(\Omega)$ is at most a singleton $\{i_{0}\}$ with $m_{t_{i_{0}}}=m(f)$. Then, we have that
\bes \left\vert f\left(\sum_{i\in\Omega}a_{i}x_{i}\right)\right\vert\leq \sum_{p=1}^{3}\left\vert f\left(\sum_{i\in A_{p}}a_{i}x_{i}\right)\right\vert + |f(a_{i_{0}}x_{i_{0}})|\leq \sum_{p=1}^{3}\left\vert f\left(\sum_{i\in A_{p}}a_{i}x_{i}\right)\right\vert + \max_{i\in\Omega}|a_{i}| \ees
because $|f(x_{i_{0}})|\leq \frac{m_{t_{i_{0}}}}{w(f)}\leq 1$. Now, recall that $\gamma=\max_{i\in\Omega}|a_{i}|$ and note that, for $A_{1}$, there is the estimate
\begin{align*} \left\vert f\left(\sum_{i\in A_{1}}a_{i}x_{i}\right)\right\vert&\leq \sum_{i\in A_{1}}|f(a_{i}x_{i})|\leq \frac{\gamma}{w(f)}\sum_{i\in A_{1}}m_{t_{i}} \\ &\leq \frac{\gamma}{w(f)}\text{card}(A_{1})m_{t_{\max(A_{1})}} \leq \frac{\gamma}{w(f)}t_{\max(A_{1})}m_{t_{\max(A_{1})}} \\ &\leq \frac{\gamma}{w(f)}(\text{ind}(m(f))-1)m_{\text{ind}(m(f))-1} \leq \frac{\gamma m_{\text{ind}(m(f))}}{w(f)\sqrt{m_{\text{ind}(m(f))}}} \leq \frac{\gamma}{\sqrt{w(f)}}\end{align*}
where we have used the first condition on the sequence $(m_{j})$ and the fact that $t_{i}<\text{ind}(m(f)):=\max_{1\leq z\leq l}j_{z}$ for each $i\in A_{1}$. Next, for $A_{2}$, note that
\begin{align*} \left\vert f\left(\sum_{i\in A_{2}}a_{i}x_{i}\right)\right\vert &\leq \frac{\gamma}{w(f)}\sum_{i\in A_{2}}m_{t_{i}}\left\vert\sum_{q=1}^{d}e^{*}_{i_{q}}\left(\sum_{r\in F_{i}}c^{i}_{r}e_{r}\right) \right\vert \\ &\leq \frac{\gamma}{w(f)}\sum_{i\in A_{2}} 3m_{t_{i}}\ep_{i} \leq\frac{\gamma}{w(f)}\sum_{i\in A_{2}}\frac{1}{2^{N}}\leq\frac{\gamma}{w(f)}\leq\frac{\gamma}{\sqrt{w(f)}} \end{align*}
where we have used the observation that immediately precedes Proposition \ref{basic-scc-maximal} and the fact that $e^{*}_{i_{1}}<\ldots<e^{*}_{i_{d}}$ are $S_{n_{j_{1}}+\cdots+n_{j_{l}}}\ast\mathcal{A}_{3}$-admissible with $n_{j_{1}}+\cdots+n_{j_{l}}<n_{t_{i}}$ for each $i\in A_{2}$. Lastly, for $A_{3}$, note that 
\begin{align*} \left\vert f\left(\sum_{i\in A_{3}}a_{i}x_{i}\right)\right\vert &\leq \frac{\gamma}{w(f)}\sum_{i\in A_{3}}m_{t_{i}} \leq \frac{\gamma}{w(f)}|A_{3}|m_{t_{\max(A_{3})}} \\ &\leq\frac{\gamma m^{2}_{t_{\max(A_{3})}}}{w(f)}\leq \frac{\gamma\sqrt{w(f)}}{w(f)}=\frac{\gamma}{\sqrt{w(f)}}\end{align*}
because $m_{t_{i}}^{2}\leq\sqrt{w(f)}$ for each $i\in A_{3}$, by Remark \ref{weight-to-admissibility condition}. It now follows by putting these pieces together that
\bes \left\vert f\left(\sum_{i\in\Omega}a_{i}x_{i}\right)\right\vert \leq \gamma+\frac{3\gamma}{\sqrt{w(f)}}\ees
for each $f\in W^{\text{aux},\delta}_{1}$ and, in particular, we have that $\left\vert f\left(\sum_{i\in A_{f}(\Omega)}a_{i}x_{i}\right)\right\vert < \frac{4\gamma}{\sqrt{w(f)}}$ as is required as part of the inductive hypothesis.

Assume now that the statement of Proposition \ref{c0_type_estimate} holds for every functional $f\in W^{\text{aux},\delta}_{k}$, let $\Omega\subset\{1,\ldots,2^{N}\}$ be given, and let 
\bes f=\frac{1}{m_{j_{1}}\cdots m_{j_{l}}}\sum_{q=1}^{d}f_{q}\in W^{\text{aux},\delta}_{k+1}\setminus W^{\text{aux},\delta}_{k} \ees
where $f_{q}\in W^{\text{aux},\delta}_{k}$ for each $1\leq q\leq d$. Define the sets $B=\{q \mid f_{q}\in W^{\text{aux},\delta}_{0}\}$ and $C=\{3,\ldots,d\}\setminus B$. We may then write $f=\sum_{\alpha=1}^{4}g_{\alpha}$ where
\bes g_{1}=\frac{1}{w(f)}f_{1},\text{ }g_{2}=\frac{1}{w(f)}f_{2},\text{ }g_{3}=\frac{1}{w(f)}\sum_{q\in B}f_{q}, \text{ and } g_{4}=\frac{1}{w(f)}\sum_{q\in C}f_{q} \ees
and it follows that
\bes \left\vert f\left(\sum_{i\in\Omega}a_{i}x_{i}\right)\right\vert \leq \left\vert f\left(\sum_{i\in A_{f}(\Omega)}a_{i}x_{i}\right)\right\vert + |f(a_{i_{0}}x_{i_{0}})| \leq \max_{i\in\Omega}|a_{i}|+\left\vert f\left(\sum_{i\in A_{f}(\Omega)}a_{i}x_{i}\right)\right\vert \ees
where, as before, $A_{f}(\Omega)=\{ i\in\Omega \mid m_{t_{i}}\neq m(f)\}$ so that $\Omega\setminus A_{f}(\Omega)$ is at most a singleton $\{i_{0}\}$ with $m_{t_{i_{0}}}=m(f)$, in which case, $|f(x_{i_{0}})|\leq \frac{m_{t_{i_{0}}}}{w(f)}\leq 1$. Let $\gamma=\max_{i\in\Omega}|a_{i}|$ as before and it now remains to estimate
\begin{align*} \left\vert f\left(\sum_{i\in A_{f}(\Omega)}a_{i}x_{i}\right)\right\vert&\leq \sum_{\alpha=1}^{4}\left\vert g_{\alpha}\left(\sum_{i\in A_{f}(\Omega)}a_{i}x_{i}\right)\right\vert \\ &\leq\frac{1}{w(f)}\left\vert f_{1}\left(\sum_{i\in A_{f}(\Omega)}a_{i}x_{i}\right)\right\vert + \frac{1}{w(f)}\left\vert f_{2}\left(\sum_{i\in A_{f}(\Omega)}a_{i}x_{i}\right)\right\vert  \\  &\qquad + \frac{1}{w(f)}\left\vert \sum_{q\in B}f_{q}\left(\sum_{i\in A_{f}(\Omega)}a_{i}x_{i}\right)\right\vert + \frac{1}{w(f)}\left\vert\sum_{q\in C}f_{q}\left(\sum_{i\in A_{f}(\Omega)}a_{i}x_{i}\right)\right\vert  \end{align*} 
where, after applying the inductive hypothesis to the $g_{1}$ and $g_{2}$ terms and using the same argument as is used in the base case for the $g_{3}$ term, this quantity is bounded above by
\begin{align*} &\frac{\gamma}{w(f)}\left(1+\frac{3}{\sqrt{w(f_{1})}}\right)+\frac{\gamma}{w(f)}\left(1+\frac{3}{\sqrt{w(f_{2})}}\right) +\frac{3\gamma}{\sqrt{w(f)}}+\frac{1}{w(f)}\left\vert \sum_{q\in C}f_{q}\left(\sum_{i\in A_{f}(\Omega)}a_{i}x_{i}\right)\right\vert \\  &\qquad\leq \frac{\frac{13}{5}\gamma}{w(f)}+\frac{3\gamma}{\sqrt{w(f)}}+\frac{1}{w(f)}\left\vert\sum_{q\in C}f_{q}\left(\sum_{i\in A_{f}(\Omega)}a_{i}x_{i}\right)\right\vert \\ &\qquad\leq \frac{(3+\frac{13}{50})\gamma}{\sqrt{w(f)}}+ \frac{1}{w(f)}\left\vert\sum_{q\in C}f_{q}\left(\sum_{i\in A_{f}(\Omega)}a_{i}x_{i}\right)\right\vert \end{align*}
because $m_{1}\geq 100$. Now, note that there are sets $E_{0}\cup E_{1}\cup\ldots\cup E_{l}$ from Lemma \ref{aux_lemma} such that $\{q \mid f_{q}\in \mathcal{C}\}=E_{0}\cup E_{1}\cup\ldots\cup E_{l}$ and such that, for each $1\leq z\leq l$, we have that $(f_{q})_{q\in E_{z}}$ is $S_{n_{j_{1}}+\cdots+n_{j_{z}}}\ast\mathcal{A}_{3}$-admissible and $\max_{q\in E_{z}}d(m(f_{q}),\lambda_{z})\leq\delta$ for some $\lambda_{z}\in\overline{\mathcal{D}}$. Next, let $\Gamma_{z}=\{ i\in A_{f}(\Omega) \mid m_{t_{i}}=m(f_{q}) \text{ for some } q\in E_{z}\}$ for each $1\leq z\leq l$ and note that, if $i_{1},i_{2}\in \Gamma_{z}$, then we have that
\begin{align*} 2^{-N} &\leq d(m_{t_{i_{1}}},m_{t_{i_{2}}}) \\ &= d(m(f_{q_{1}}),m(f_{q_{2}}))\leq d(m(f_{q_{1}}),\lambda_{z})+d(\lambda_{z},m(f_{q_{2}})) \leq 2\delta \end{align*}
for the corresponding $q_{1},q_{2}\in E_{z}$, and this cannot happen. It follows that $\Gamma_{z}$ is at most a singleton, say $i_{z}$, for each $1\leq z\leq l$ so that
\begin{align*}\left\vert \sum_{q\in\mathcal{C}}f_{q}\left(\sum_{i\in A_{f}(\Omega)}a_{i}x_{i}\right)\right\vert &\leq \sum_{z=0}^{l} \sum_{q\in E_{z}} \left\vert f_{q}\left(\sum_{i \in A_{f}(\Omega)\setminus\{i_{z}\}}a_{i}x_{i}\right)\right\vert + \sum_{z=0}^{l}\sum_{q\in E_{z}}|f_{q}(a_{i_{z}}x_{i_{z}})| \\ &\leq \sum_{z=0}^{l}\sum_{q\in E_{z}}\frac{4\gamma}{\sqrt{w(f_{q})}} + \gamma\sum_{z=0}^{l}\sum_{q\in E_{z}}\frac{m_{t_{i_{z}}}}{w(f_{q})} \\ &\leq (l+1)\gamma + \sum_{z=0}^{l}\sum_{q\in E_{z}}\frac{4\gamma}{\sqrt{w(f_{q})}} \leq (l+5)\gamma  \end{align*}
where again we have used the inductive hypothesis and, by putting all of these pieces together, there is the estimate
\begin{align*} \left\vert f\left(\sum_{i\in\Omega}a_{i}x_{i}\right)\right\vert &\leq \gamma+\frac{(3+\frac{13}{50})\gamma}{\sqrt{w(f)}}+\frac{(l+5)\gamma}{w(f)} \\ &\leq \gamma+\frac{(3+\frac{13}{50})\gamma}{\sqrt{w(f)}}+\frac{(l+5)\gamma}{10^{l}\sqrt{w(f)}} \leq \gamma + \frac{(3+\frac{13}{50}+\frac{3}{5})\gamma}{\sqrt{w(f)}} \leq \gamma+ \frac{4\gamma}{\sqrt{w(f)}}\end{align*}
where, in particular, $\left\vert f\left(\sum_{i\in A_{f}(\Omega)}a_{i}x_{i}\right)\right\vert\leq \frac{4\gamma}{\sqrt{w(f)}} $ so that the proof is complete.
\end{proof}

\end{document}
